%% macros/figure_macros.tex -- the figure macro family.
%%
%% The whole family lives in this one file.  Float placement comes from
%% \thesisFloatSpec ([template] float_spec in thesis.toml), and the
%% r-suffixed macros below always clip what they trim.
%%
%% Naming scheme
%% -------------
%%   \fig      single figure, floating
%%   \hfig     the same, placed exactly here (no float)
%%   \mpfig    one panel of a multi-panel figure (inside a figure environment)
%%
%% Suffixes compose:  t = extra text under the caption
%%                    r = trim and clip the graphic
%%                    a = explicit vertical alignment for a panel
%%
%% Every macro takes the *width* first and the *file* second, except the
%% alignment-aware panel macros where the file is third.  thesis-cli has
%% the same table built in and uses it to find every graphic the document
%% references.  A macro you add here also goes under [figures.macros] in
%% thesis.toml (macro name = position of the file argument), or its
%% graphics are missed.

\makeatletter

%% Default float placement. Overridden from thesis.toml via
%% config/generated.tex ([template] float_spec).
\newcommand{\thesisFloatSpec}{htbp}

%% Opening a float with a placement specifier that lives in a macro needs
%% one expansion step: LaTeX's float code walks the specifier token by token
%% and would otherwise see the macro itself rather than "htbp".
\newcommand{\ThesisBeginFloat}[1]{%
  \edef\thesis@floatbegin{\noexpand\begin{#1}[\thesisFloatSpec]}%
  \thesis@floatbegin}

%% Non-floating figures and tables.  \@captype is set by hand so that
%% \caption works outside a float environment.
\newenvironment{figurehere}{\def\@captype{figure}\par\centering}{\par}
\newenvironment{tablehere}{\def\@captype{table}\par\centering}{\par}

%% ------------------------------------------------------------------
%% Single figures, floating
%% ------------------------------------------------------------------
%% \fig{width}{file}{caption}{label}
%%   width is a fraction of \columnwidth, e.g. 0.8
\newcommand{\fig}[4]{%
  \ThesisBeginFloat{figure}%
    \centering
    \includegraphics[width=#1\columnwidth]{#2}%
    \caption{#3}\label{#4}%
  \end{figure}}

%% \figt{width}{file}{caption}{label}{extra text below the caption}
\newcommand{\figt}[5]{%
  \ThesisBeginFloat{figure}%
    \centering
    \includegraphics[width=#1\columnwidth]{#2}%
    \caption{#3}\label{#4}%
    \par\medskip
    #5%
  \end{figure}}

%% \figtr{width}{file}{caption}{label}{trim}
%%   trim is "left bottom right top", unitless values are big points (bp)
\newcommand{\figtr}[5]{%
  \ThesisBeginFloat{figure}%
    \centering
    \includegraphics[width=#1\columnwidth,trim={#5},clip=true]{#2}%
    \caption{#3}\label{#4}%
  \end{figure}}

%% ------------------------------------------------------------------
%% Single figures, placed exactly here
%% ------------------------------------------------------------------
\newcommand{\hfig}[4]{%
  \begin{figurehere}%
    \includegraphics[width=#1\columnwidth]{#2}%
    \caption{#3}\label{#4}%
  \end{figurehere}}

\newcommand{\hfigt}[5]{%
  \begin{figurehere}%
    \includegraphics[width=#1\columnwidth]{#2}%
    \caption{#3}\label{#4}%
    \par\medskip
    #5%
  \end{figurehere}}

\newcommand{\hfigtr}[5]{%
  \begin{figurehere}%
    \includegraphics[width=#1\columnwidth,trim={#5},clip=true]{#2}%
    \caption{#3}\label{#4}%
  \end{figurehere}}

%% ------------------------------------------------------------------
%% Panels of a multi-panel figure
%% ------------------------------------------------------------------
%% These go *inside* a figure environment, side by side:
%%
%%   \begin{figure}[htbp]\centering
%%     \mpfig{0.45}{figures/left}{Left panel}{fig:left}\hfill
%%     \mpfig{0.45}{figures/right}{Right panel}{fig:right}
%%   \end{figure}
%%
%% The minipage is given a height of 1\totalheight so that panels of
%% different heights still align on their bottom edge.

%% \mpfig{width}{file}{caption}{label}
\newcommand{\mpfig}[4]{%
  \begin{minipage}[b][1\totalheight]{#1\columnwidth}%
    \centering
    \includegraphics[width=\columnwidth]{#2}%
    \caption{#3}\label{#4}%
  \end{minipage}}

%% \mpfiga{width}{alignment}{file}{caption}{label}
%%   alignment is b, c or t
\newcommand{\mpfiga}[5]{%
  \begin{minipage}[#2][1\totalheight]{#1\columnwidth}%
    \centering
    \includegraphics[width=\columnwidth]{#3}%
    \caption{#4}\label{#5}%
  \end{minipage}}

%% \mpfigt{width}{file}{caption}{label}{extra text}
\newcommand{\mpfigt}[5]{%
  \begin{minipage}[b][1\totalheight]{#1\columnwidth}%
    \centering
    \includegraphics[width=\columnwidth]{#2}%
    \caption{#3}\label{#4}%
    \par\medskip
    #5%
  \end{minipage}}

%% \mpfigtr{width}{file}{caption}{label}{trim}
\newcommand{\mpfigtr}[5]{%
  \begin{minipage}[b][1\totalheight]{#1\columnwidth}%
    \centering
    \includegraphics[width=\columnwidth,trim={#5},clip=true]{#2}%
    \caption{#3}\label{#4}%
  \end{minipage}}

%% \mpfigtra{width}{alignment}{file}{caption}{label}{trim}
\newcommand{\mpfigtra}[6]{%
  \begin{minipage}[#2][1\totalheight]{#1\columnwidth}%
    \centering
    \includegraphics[width=\columnwidth,trim={#6},clip=true]{#3}%
    \caption{#4}\label{#5}%
  \end{minipage}}

%% \mpfigabc{width}{file}
%%   A panel with no caption of its own, for (a)(b)(c) sub-figures that
%%   share one caption placed after the last panel.
\newcommand{\mpfigabc}[2]{%
  \begin{minipage}[b][1\totalheight]{#1\columnwidth}%
    \centering
    \includegraphics[width=\columnwidth]{#2}%
  \end{minipage}}

\makeatother
